% Part: many-valued-logic
% Chapter: sequent-calculus
% Section: structural-rules

\documentclass[../../../include/open-logic-section]{subfiles}

\begin{document}

\olfileid{mvl}{seq}{str}

\olsection{నిర్మాణాత్మక నియమాలు}

$n$-వైపుల సీక్వెంట్ కలనంలోని
నిర్మాణాత్మక నియమాలు సాంప్రదాయిక
సందర్భంలో మాదిరిగానే పనిచేస్తాయి;
అయితే ఇక్కడ ప్రతి స్థానం~$i$కూ
అవి వర్తిస్తాయి.

\begin{defish}
\begin{center}
\AxiomC{$ \Gamma_1 \nSequent \dots \nSequent \phantom{!A,}\Gamma_i \nSequent \dots \nSequent \Gamma_n $}
\RightLabel{$\iR{\Weakening}{i}$}
\UnaryInfC{$\Gamma_1 \nSequent \dots \nSequent !A, \Gamma_i \nSequent \dots \nSequent \Gamma_n$}
\DisplayProof
\\[2ex]
\AxiomC{$ \Gamma_1 \nSequent \dots \nSequent !A, !A, \Gamma_i \nSequent \dots \nSequent \Gamma_n $}
\RightLabel{$\iR{\Contraction}{i}$}
\UnaryInfC{$\Gamma_1 \nSequent \dots \nSequent \phantom{!A,}!A, \Gamma_i \nSequent
\dots \nSequent \Gamma_n$}
\DisplayProof
\\[2ex]
\AxiomC{$ \Gamma_1 \nSequent \dots \nSequent \Gamma_i, !A, !B, \Gamma_i' \nSequent \dots \nSequent \Gamma_n $}
\RightLabel{$\iR{\Exchange}{i}$}
\UnaryInfC{$\Gamma_1 \nSequent \dots \nSequent \Gamma_i, !B, !A, \Gamma_i' \nSequent \dots
\nSequent \Gamma_n$}
\DisplayProof
\end{center}
\end{defish}

బలహీనీకరణ, సంకోచనం, మార్పిడి
నిగమనాల వరుసను తరచుగా
ద్విరేఖ నిగమన గీతలతో సూచిస్తాం.

\Cut{} నియమానికి అనేక రూపాలు ఉన్నాయి.
సీక్వెంట్‌లోని వేర్వేరు స్థానాలు
$i \neq j$ జతలలో ప్రతిదానికీ
ఒక రూపం ఉంటుంది:
\begin{defish}
\[
\AxiomC{$ \Gamma_1 \nSequent \dots \nSequent !A, \Gamma_i \nSequent \dots \nSequent \Gamma_n $}
\AxiomC{$ \Delta_1 \nSequent \dots \nSequent !A, \Delta_j \nSequent \dots \nSequent \Delta_n $}
\RightLabel{$\iR{\Cut}{i,j}$}
\BinaryInfC{$\Gamma_1,\Delta_1 \nSequent \dots \nSequent \Gamma_n, \Delta_n$}
\DisplayProof
\]
\end{defish}

\end{document}
